\subsection{The complete channel table}\label{sec:superspace:table}

We now state all seven nonzero families of the master
identity~\eqref{eq:ss:master}. The letter
definitions~\eqref{eq:ss:letters} induce a uniform rescaling between the
unit-normalized field-strength composites used in the graph computations and
the letters of Definition~\ref{def:ss:letters}. The bookkeeping is immediate:
in the seed channel the two insertions contribute $(-i/\sqrt2)^2=-\tfrac12$
and the output bilinears contribute $\sqrt2$, giving the net prefactor
$-\tfrac12\cdot\sqrt2\cdot\tfrac{\h g^2}{16\pi^2}
=-\tfrac{\h g^{2}}{16\sqrt2\,\pi^{2}}$ of equation~\eqref{eq:ss:bbresult}.
The same arithmetic, applied family by family, yields the common prefactor
$-\h g^{2}/(16\sqrt2\,\pi^{2})$ for every family, including the
$(b,\beta)/(\beta,b)$ family, whose output word is fixed by the
residual-supersymmetry closure discussed below.
With $\langle X,Y\rangle$ as in~\eqref{eq:ss:bilinear} and
$\partial_{\dot\alpha}(\partial c)$ the holomorphic derivative of the letter
$\partial c$, the complete set of nonzero channels is
\begin{align}
\nabla_{-}\big(b^{A}b^{B}\big)
&=-\frac{\h g^{2}}{16\sqrt2\,\pi^{2}}\,f^{ACD}f^{BCE}\,
\Big\{\langle\partial c,b\rangle-\langle b,\partial c\rangle
+\textstyle\sum_r\big(\langle\beta_r,\gamma^r\rangle
-\langle\gamma^r,\beta_r\rangle\big)\Big\}^{DE},
\label{eq:ss:ch-bb}\\[4pt]
\nabla_{-}\big(b^{A}\gamma^{rB}\big),\quad
\nabla_{-}\big(\gamma^{rA}b^{B}\big)
&=-\frac{\h g^{2}}{16\sqrt2\,\pi^{2}}\,f^{ACD}f^{BCE}\,
\big\{\langle\partial c,\gamma^{r}\rangle
-\langle\gamma^{r},\partial c\rangle\big\}^{DE},
\label{eq:ss:ch-bgamma}\\[4pt]
\nabla_{-}\big(\beta_{r}^{A}\gamma^{sB}\big),\quad
\nabla_{-}\big(\gamma^{sA}\beta_{r}^{B}\big)
&=-\frac{\h g^{2}}{16\sqrt2\,\pi^{2}}\,f^{ACD}f^{BCE}\,
\delta_{rs}\,\langle\partial c,\partial c\rangle^{DE},
\label{eq:ss:ch-betagamma}\\[4pt]
\nabla_{-}\big(\beta_{r}^{A}\beta_{s}^{B}\big)
&=-\frac{\h g^{2}}{16\sqrt2\,\pi^{2}}\,f^{ACD}f^{BCE}\,
\varepsilon_{rst}\,
\big\{\langle\partial c,\gamma^{t}\rangle
-\langle\gamma^{t},\partial c\rangle\big\}^{DE},
\label{eq:ss:ch-betabeta}\\[4pt]
\nabla_{-}\big(b^{A}\beta_{r}^{B}\big),\quad
\nabla_{-}\big(\beta_{r}^{A}b^{B}\big)
&=-\frac{\h g^{2}}{16\sqrt2\,\pi^{2}}\,f^{ACD}f^{BCE}\,
\big\{\langle\beta_{r},\partial c\rangle
+\langle\partial c,\beta_{r}\rangle
+\varepsilon_{rst}\langle\gamma^{s},\gamma^{t}\rangle\big\}^{DE},
\label{eq:ss:ch-bbeta}\\[4pt]
\nabla_{-}\big(b^{A}\partial_{\dot\alpha}c^{B}\big)
&=-\frac{\h g^{2}}{16\sqrt2\,\pi^{2}}\,f^{ACD}f^{BCE}\,
\Big\{\tfrac13\big\langle\partial_{\dot\alpha}(\partial c),\partial c\big\rangle
+\tfrac23\big\langle\partial c,\partial_{\dot\alpha}(\partial c)\big\rangle
\Big\}^{DE},
\label{eq:ss:ch-bdc}\\[4pt]
\nabla_{-}\big(\partial_{\dot\alpha}c^{A}b^{B}\big)
&=-\frac{\h g^{2}}{16\sqrt2\,\pi^{2}}\,f^{ACD}f^{BCE}\,
\Big\{\tfrac23\big\langle\partial_{\dot\alpha}(\partial c),\partial c\big\rangle
+\tfrac13\big\langle\partial c,\partial_{\dot\alpha}(\partial c)\big\rangle
\Big\}^{DE}.
\label{eq:ss:ch-dcb}
\end{align}
The free dotted index in~\eqref{eq:ss:ch-bdc} and~\eqref{eq:ss:ch-dcb} is the
letter index of $\partial_{\dot\alpha}c$ on the left and the extra holomorphic
derivative on the right; the weights $(\tfrac13,\tfrac23)$ are the
$(m,n)=(1,0)$ instance of the tower kernel of
Section~\ref{sec:superspace:tower}. Together with the fifty-two zeros of
Section~\ref{sec:superspace:selection}, the rows
\eqref{eq:ss:ch-bb}--\eqref{eq:ss:ch-dcb} exhaust the eighty-one ordered
channels: $1+6+6+4+6+6=29$ nonzero, as in~\eqref{eq:ss:29}.

\subsubsection*{Residual-supersymmetry closure of the $(b,\beta)/(\beta,b)$
channels}

The family~\eqref{eq:ss:ch-bbeta} is the only one whose output word is not
read off a cut parent graph directly. The raw graph carrier of this family
contains an equation-of-motion (total-derivative) layer; the physical answer
is the unique residual-supersymmetry-closed vector obtained after removing
that layer and adding a finite normal-product (scheme) term, with the
remaining one-parameter freedom anchored to the $(b,b)$ scale. The closure
works as follows. The residual supercharges of the bottom-letter sector,
$q_r:=\tfrac{1}{\sqrt2}Q^{r}_{E,+}$, act on the letters by
\begin{equation}\label{eq:ss:qaction}
q_r\,b=0,
\qquad
q_r\,\beta_s=\delta_{rs}\,b,
\qquad
q_r\,\gamma^{s}=-\varepsilon_{rst}\,\beta_t,
\qquad
q_r\,(\partial_{\dot\alpha}c)=\partial_{\dot\alpha}\gamma^{r} .
\end{equation}
At twisted dimension $\tfrac92$, the space of odd $SU(3)$-singlet letter
bilinears annihilated by all $q_r$ is one-dimensional, spanned by the seed
output word
\begin{equation}\label{eq:ss:singlet}
\big\{\langle\partial c,b\rangle-\langle b,\partial c\rangle
+\textstyle\sum_r\big(\langle\beta_r,\gamma^r\rangle
-\langle\gamma^r,\beta_r\rangle\big)\big\}^{DE}
\end{equation}
of~\eqref{eq:ss:ch-bb}. The $(b,\beta_r)$ output word is its unique triplet
$q$-primitive:
\begin{equation}\label{eq:ss:triplet}
q_s\big\{\langle\beta_r,\partial c\rangle+\langle\partial c,\beta_r\rangle
+\varepsilon_{rtu}\langle\gamma^{t},\gamma^{u}\rangle\big\}
=-\delta_{sr}\,
\Big\{\langle\partial c,b\rangle-\langle b,\partial c\rangle
+\textstyle\sum_{r'}\big(\langle\beta_{r'},\gamma^{r'}\rangle
-\langle\gamma^{r'},\beta_{r'}\rangle\big)\Big\} .
\end{equation}
Equivalently, in the ordered basis
$\big(\langle\partial c,\beta_r\rangle,\,\langle\beta_r,\partial c\rangle,\,
\langle\gamma^{s},\gamma^{t}\rangle,\,\langle\gamma^{t},\gamma^{s}\rangle\big)$
with $\varepsilon_{rst}=+1$, the three Ward equations $q_s[\,\cdot\,]=0$
($s=1,2,3$) cut out the one-dimensional ray spanned by $(1,1,1,-1)$, and
anchoring this ray to the $(b,b)$ scale through~\eqref{eq:ss:triplet} fixes
the unit normalization displayed in~\eqref{eq:ss:ch-bbeta}. Since
$q_r(b^Ab^B)=0$, the seed channel's own forward orbit is only itself: the
physical families were each computed directly, and the closure is a
cross-check that fixes the $(b,\beta)$ normalization.

\begin{remark}[The finite shift is a scheme term, not an anomaly graph]%
\label{rem:ss:finiteshift}
The passage from the raw graph carrier to the renormalized coefficient vector
of the $(b,\beta)/(\beta,b)$ family removes an equation-of-motion
(total-derivative) layer and adds a finite normal-product term. This is a
finite scheme adjustment, not an additional anomaly graph: all bare-graph
anomaly in every channel comes only from the $\mul^{2}$ cutting
failure~\eqref{eq:ss:cuttingfailure}. The closure was verified by thirty-one
independent finite Ward-identity checks. Through this finite shift and the
anchoring to the $(b,b)$ scale, the $(b,\beta)/(\beta,b)$ family lands on the
same universal prefactor $-\h g^{2}/(16\sqrt2\,\pi^{2})$ as every other
nonzero family, with the output word displayed in~\eqref{eq:ss:ch-bbeta}.
\end{remark}

\subsection{The holomorphic-derivative tower}\label{sec:superspace:tower}

Shifting the second letter of a channel bilinear by a formal vector $w$ ---
equivalently, resumming all holomorphic derivatives --- inserts the leg phase
$e^{iw\cdot(aq+bp)}$ into the ordered Feynman parametrization of the triangle.
The ordered-simplex moment is
\begin{equation}\label{eq:ss:simplexmoment}
\int_{0<a<b<1}a^{p}b^{q}\,da\,db=\frac{1}{(p+1)(p+q+2)} ,
\end{equation}
so that
\begin{equation}\label{eq:ss:simplexderivation}
2\int_{0}^{1}\!db\int_{0}^{b}\!da\;a^{k+\ell}\,b^{(m-k)+(n-\ell)}
=\frac{2}{(k+\ell+1)(m+n+2)} .
\end{equation}
Restoring the binomial weights of the derivative expansion gives the one-loop
holomorphic-derivative kernel
\begin{equation}\label{eq:ss:kernel}
\boxed{\;
K_{m,n;k,\ell}
=\frac{2\binom{m}{k}\binom{n}{\ell}}{(m+n+2)(k+\ell+1)}\;}
\qquad
(m,n\in\mathbb Z_{\geq0},\ 0\leq k\leq m,\ 0\leq\ell\leq n).
\end{equation}

\begin{remark}[The overall factor $2$ is $\Gamma(3)$]\label{rem:ss:gamma3}
The numerator factor $2$ in~\eqref{eq:ss:kernel} is
$\Gamma(3)$ from combining the three denominators,
$\frac{1}{D_0D_1D_2}=2\int_{\Delta_2}\frac{1}{(r^{2}+\Delta)^{3}}$; it is not
an orientation multiplicity and not a second loop orientation.
\end{remark}

Because the $\partial_{\dot\alpha}$ are gauge-covariant and do not commute,
derivative monomials are understood as shuffle-symmetrized averages over the
$(m,n)$ orderings of $\partial_{\dot1}^{m}\partial_{\dot2}^{n}$. The full
kernel operator on letter pairs is
\begin{equation}\label{eq:ss:kerneloperator}
\mathcal K_{m,n}(f^{D},g^{E})
=\sum_{k=0}^{m}\sum_{\ell=0}^{n}
\frac{2\binom{m}{k}\binom{n}{\ell}}{(m+n+2)(k+\ell+1)}\,
\big(\partial_{\dot1}^{k}\partial_{\dot2}^{\ell}\partial_{\dot\alpha}f\big)^{D}
\big(\partial_{\dot1}^{\,m-k}\partial_{\dot2}^{\,n-\ell}
\partial^{\dot\alpha}g\big)^{E} .
\end{equation}
Its first instances are
\begin{align}
\mathcal K_{0,0}&=\langle f,g\rangle,
\qquad
\mathcal K_{1,0}=\frac13\langle\partial_{\dot1}f,g\rangle
+\frac23\langle f,\partial_{\dot1}g\rangle,
\label{eq:ss:kernelexamples1}\\
\mathcal K_{2,0}&=\frac16\langle\partial_{\dot1}^{2}f,g\rangle
+\frac12\langle\partial_{\dot1}f,\partial_{\dot1}g\rangle
+\frac12\langle f,\partial_{\dot1}^{2}g\rangle,
\label{eq:ss:kernelexamples2}\\
\mathcal K_{1,1}&=\frac16\langle\partial_{\dot1}\partial_{\dot2}f,g\rangle
+\frac14\langle\partial_{\dot1}f,\partial_{\dot2}g\rangle
+\frac14\langle\partial_{\dot2}f,\partial_{\dot1}g\rangle
+\frac12\langle f,\partial_{\dot1}\partial_{\dot2}g\rangle .
\label{eq:ss:kernelexamples3}
\end{align}
In particular the $(b,\partial c)$ weights $(\tfrac13,\tfrac23)$ of
\eqref{eq:ss:ch-bdc}--\eqref{eq:ss:ch-dcb} are the $(m,n)=(1,0)$ instance:
$K_{1,0;1,0}=\tfrac{2}{3\cdot2}=\tfrac13$ for the derivative on the first
letter and $K_{1,0;0,0}=\tfrac{2}{3\cdot1}=\tfrac23$ for the derivative on
the second. The comparison of~\eqref{eq:ss:kernel} with the holomorphic-twist
kernel of \cite{BK} is the subject of Section~\ref{sec:ht}.
